\documentclass[11pt,a4paper]{article}

\usepackage[utf8]{inputenc}
\usepackage[T1]{fontenc}
\usepackage[margin=1in]{geometry}
\usepackage{graphicx}
\usepackage{amsmath}
\usepackage{booktabs}
\usepackage{listings}
\usepackage{xcolor}
\usepackage{enumitem}
\usepackage{hyperref}
\hypersetup{colorlinks=true, linkcolor=blue!70!black, urlcolor=blue!70!black,
            citecolor=blue!70!black, pdftitle={Lab No. 07 -- Jenkins CI/CD Pipeline}}

% ---------------------------------------------------------------------------
% TO EDIT: put your own username in the repository URL below.
% (Your previous lab report used the GitHub username "RebahiWassim".)
% ---------------------------------------------------------------------------
\newcommand{\studentName}{Rebahi M'hamed Wassim}
\newcommand{\studentGroup}{2}
\newcommand{\repoURL}{https://github.com/RebahiWassim/jenkins-flask-lab.git}
% ---------------------------------------------------------------------------

\lstdefinestyle{code}{
  basicstyle=\ttfamily\small,
  breaklines=true,
  frame=single,
  framesep=5pt,
  xleftmargin=4pt,
  backgroundcolor=\color{gray!8},
  commentstyle=\color{gray},
  keywordstyle=\color{blue!70!black},
  stringstyle=\color{red!70!black}
}
\lstset{style=code}

% Inserts the screenshot if the file exists; otherwise shows a labelled box.
\newcommand{\evidence}[2]{%
  \begin{figure}[htbp]
    \centering
    \IfFileExists{#2}{%
      \includegraphics[width=0.88\textwidth]{#2}%
    }{%
      \fbox{\parbox[t]{0.84\textwidth}{\centering\vspace*{0.7cm}
        \textbf{Screenshot to insert: #1}\\[1mm]
        {\small Save the screenshot as \texttt{#2} in this folder and recompile;
        it will be inserted automatically.}%
        \vspace*{0.7cm}}%
      }%
    }
    \caption{#1}
  \end{figure}
}

\begin{document}

\begin{center}
{\Large\bfseries Lab No. 07 --- Building a CI/CD Pipeline with Jenkins and Docker}\\[3mm]
Student: \studentName \hfill Group: \studentGroup\\
Level: 5th-year Software Engineering \quad | \quad Course: DevOps \& Cloud Computing\\
Lab date: 22/11/2026 \quad | \quad Report: \today\\[2mm]
GitHub repository: \href{\repoURL}{\texttt{\repoURL}}
\end{center}

\hrule\vspace{3mm}

\tableofcontents\vspace{2mm}

\hrule\vspace{4mm}

\section{Introduction and objectives}

This lab puts into practice a complete Continuous Integration / Continuous Delivery
workflow. A small Flask application is stored on GitHub, built and tested by Jenkins
inside Docker images, then deployed as a container and checked over HTTP. The goal is
to understand how an automated pipeline turns a commit into a verified running service,
and why a failing test must block deployment.

The workflow is:

\[
  \text{GitHub} \;\longrightarrow\; \text{Jenkins Pipeline} \;\longrightarrow\;
  \text{Docker Image} \;\longrightarrow\; \text{Docker Container} \;\longrightarrow\;
  \text{Flask Application}
\]

Jenkins is reached at \texttt{http://localhost:8080} and the deployed application at
\texttt{http://localhost:8081}. Jenkins runs on an agent that has access to the Docker
daemon; this is acceptable only on the isolated educational machine used for the lab,
never on a shared or production host.

\section{Project structure and application}

The repository contains the following files:

\begin{lstlisting}[basicstyle=\ttfamily\footnotesize]
jenkins-flask-lab/
 |- app.py               Flask app: "/" (home) and "/health" endpoints
 |- requirements.txt     Flask, pytest, gunicorn
 |- Dockerfile           python:3.12-slim image, gunicorn on port 8000
 |- Jenkinsfile          declarative 5-stage pipeline
 |- tests/test_app.py    pytest tests for both endpoints
\end{lstlisting}

\texttt{app.py} reads the \texttt{APP\_VERSION} environment variable and displays it on
the home page; the \texttt{/health} endpoint returns \texttt{\{"status": "ok"\}} with
HTTP 200. The tests use Flask's test client to check both endpoints. The Dockerfile
copies the application {\em and} its tests into the image, so the very same image can be
used to run the tests and to serve the application with gunicorn.

\section{The Jenkins pipeline}

The \texttt{Jenkinsfile} is stored at the root of the repository and defines a
declarative pipeline with three environment variables: \texttt{IMAGE\_NAME}
(\texttt{jenkins-flask-lab}), \texttt{CONTAINER\_NAME} (\texttt{jenkins-flask-app})
and \texttt{APP\_PORT} (\texttt{8081}). Every image is tagged with
\texttt{\$\{IMAGE\_NAME\}:\$\{BUILD\_NUMBER\}}, which keeps each build traceable;
the build number is also injected as \texttt{APP\_VERSION} into the running container.

\begin{table}[htbp]
\centering
\begin{tabular}{@{}p{0.16\textwidth}p{0.42\textwidth}p{0.34\textwidth}@{}}
\toprule
\textbf{Stage} & \textbf{Command / action} & \textbf{Expected result}\\
\midrule
Checkout & Jenkins clones the SCM repository into the workspace (``Pipeline script from SCM''); the stage only echoes a message to make the step visible. & Source code and Jenkinsfile are available.\\
Build & \texttt{docker build -t \$\{IMAGE\_NAME\}:\$\{BUILD\_NUMBER\} .} & A tagged Docker image is created.\\
Test & \texttt{docker run --rm \$\{IMAGE\_NAME\}:\$\{BUILD\_NUMBER\} pytest -v} & Tests pass; the temporary container is removed. A failure stops the pipeline.\\
Deploy & \texttt{docker rm -f \$\{CONTAINER\_NAME\} || true}, then \texttt{docker run -d --name \$\{CONTAINER\_NAME\} -p \$\{APP\_PORT\}:8000 -e APP\_VERSION=\$\{BUILD\_NUMBER\} \$\{IMAGE\_NAME\}:\$\{BUILD\_NUMBER\}} & The previous container is replaced by the new version on port 8081.\\
Verify & \texttt{curl --fail --retry 10 --retry-delay 2 http://localhost:\$\{APP\_PORT\}/health} and \texttt{curl --fail http://localhost:\$\{APP\_PORT\}/} & Health endpoint and home page respond successfully.\\
\bottomrule
\end{tabular}
\caption{The five stages of the pipeline and their expected results.}
\end{table}

The \texttt{post} block prints a message on success or failure. The \texttt{--rm} flag
prevents test containers from accumulating, \texttt{|| true} makes the Deploy stage
idempotent when no previous container exists, and the retry options give gunicorn time
to start before the health check is considered failed.

\section{Running and verifying the pipeline}

A Jenkins Pipeline job named \texttt{jenkins-flask-lab} was created with
\textit{Pipeline script from SCM}, Git as SCM, the public repository URL, branch
\texttt{*/main} and script path \texttt{Jenkinsfile}. After \textit{Build Now}, the
console output shows the five stages executing in order and build \#1 ends with
\texttt{SUCCESS}. No credentials were needed because the repository is public. Two
intermediate builds were interrupted by transient GitHub connectivity during SCM
checkout; after restarting Jenkins with a short DNS cache TTL, the final build \#5
also completed successfully and deployed version 5 of the application.

\evidence{Jenkins Pipeline Steps view for build \#5: all five stages succeeded.}{jenkins-success.png}

The application was then checked on the Jenkins host:

\begin{lstlisting}
curl http://localhost:8081/health   # {"status":"ok"}
curl http://localhost:8081/         # <h1>Jenkins CI/CD Lab</h1><p>Version: <BUILD_NUMBER></p>
\end{lstlisting}

\evidence{Flask application running in the browser at \texttt{http://localhost:8081} with ``Version: 5''.}{app-running.png}

\section{Failure experiment}

To show that testing gates delivery, the health assertion in \texttt{tests/test\_app.py}
was temporarily changed to \texttt{assert response.get\_json() == \{"status":
"incorrect"\}} and pushed. Build \#2 failed in the \textbf{Test} stage: pytest
reported \texttt{AssertionError: assert \{'status': 'ok'\} == \{'status':
'incorrect'\}}, so Jenkins stopped the pipeline. The \textbf{Deploy} and
\textbf{Verify} stages were skipped, the \texttt{post failure} message was printed
and the build was marked \texttt{FAILURE}.
Consequently, the previously running container was not replaced by the broken version ---
the failed code never reached the users. Restoring \texttt{\{"status": "ok"\}}, pushing
the fix and rebuilding produced a successful pipeline again.

\evidence{Failed Test stage in the Jenkins console output (build \#2).}{test-failed.png}

\evidence{Jenkins Pipeline Steps view of build \#2: Test failed and Deploy/Verify were skipped.}{jenkins-failed-stages.png}

\section{Blue--Green deployment and rollback (Section 19)}

\begin{enumerate}[label=\arabic*.]
\item \textbf{Deployment vs.\ release.} Deployment is placing a new version in an
  environment (building the image, starting the containers, making it reachable there).
  Release is routing real user traffic to that version. With Blue--Green, Green can be
  deployed and validated while Blue still serves all production traffic; the release
  happens only when traffic is switched.
\item \textbf{Why deploy to Green first?} It provides a production-like environment for
  final tests without disturbing current users. Errors, performance regressions or a
  bad configuration are detected on Green while Blue continues to serve traffic, so
  the impact of a faulty release is close to zero.
\item \textbf{Which component routes traffic?} A reverse proxy or load balancer such as
  NGINX, HAProxy or a cloud load balancer (a DNS switch or service mesh can also be
  used). It can point user requests to the Blue or the Green environment by changing
  its upstream configuration.
\item \textbf{Why keep Blue available?} It is the known-good rollback target. If Green
  fails after the switch, traffic can be sent back to Blue immediately, without
  rebuilding or redeploying the old version.
\item \textbf{How can Jenkins automate it?} Jenkins can build and tag images, run unit
  tests, security scans and integration/smoke tests against Green, deploy Green, run
  health checks, wait for a manual approval with an \texttt{input} step, then trigger the
  proxy reconfiguration (script, API or infrastructure-as-code), and finally monitor the
  result and roll back automatically if checks fail. Credentials are supplied through
  Jenkins credentials bindings, never the Jenkinsfile.
\item \textbf{Main benefit of rollback.} Fast recovery with minimal downtime: returning
  to the previous working version is just a traffic switch, not a new build and
  deployment, which greatly reduces risk and user-visible incidents.
\end{enumerate}

\section{Questions (Section 20)}

\begin{enumerate}[label=\arabic*.]
\item \textbf{What is Jenkins and what problem does it solve?} Jenkins is an open-source
  automation server that orchestrates repeatable jobs through pipelines. It solves the
  manual, error-prone and inconsistent build/test/deploy process by automatically
  fetching code, building it, running tests and deploying it, with a visible and
  auditable history for every build.
\item \textbf{Purpose of a Jenkinsfile.} It is ``pipeline as code'': a text file (Groovy
  DSL) kept in the repository next to the source. It declares the agent, environment,
  stages and post-build actions, so the pipeline is versioned, reviewable and
  reproducible instead of being configured by hand in the Jenkins UI.
\item \textbf{CI vs.\ CD.} Continuous Integration means developers merge changes
  frequently into a shared branch and every merge triggers an automated build and
  tests, catching integration problems early. Continuous Delivery extends this by
  automatically packaging and deploying the validated code to staging/production-like
  environments so that every good build is ready to be released, often with a manual
  approval before the final release. Continuous Deployment goes one step further and
  releases to production automatically.
\item \textbf{Role of GitHub here.} GitHub hosts the application source and the
  Jenkinsfile as the single source of truth. Jenkins is configured as ``Pipeline script
  from SCM'' and clones the \texttt{main} branch from the repository URL; a webhook
  could also notify Jenkins on each push (used only if Jenkins is reachable from
  GitHub).
\item \textbf{Why Docker?} It provides an isolated, reproducible environment that
  behaves the same on the developer machine and on the Jenkins agent. The same image is
  used for tests and for production, so ``it works on my machine'' issues disappear.
  Containers are fast to create and destroy, and the build-number tag makes each
  version traceable.
\item \textbf{Why test before deployment?} Tests are a quality gate: they detect
  regressions while the change is still cheap to discard. Because the pipeline stops
  at the failed stage, a broken version is never deployed and users keep seeing the
  last working release.
\item \textbf{What happens when a stage fails?} The pipeline halts at that stage; all
  later stages (Deploy and Verify in this lab) are skipped, the \texttt{post failure}
  block runs, and Jenkins marks the build as \texttt{FAILURE}. No new container is
  started, so the previous deployment remains in place.
\item \textbf{Purpose of the Verify stage.} Starting a container only proves that Docker
  accepted the command. Verify sends real HTTP requests: the health endpoint is retried
  for up to 20 seconds to absorb startup time, and the home page is fetched with
  \texttt{curl --fail}. The build succeeds only when the deployed service actually
  answers.
\item \textbf{Why not store credentials in a Jenkinsfile?} The Jenkinsfile lives in the
  repository and is readable by everyone with repository access and by anyone viewing
  build logs; secrets would leak through history and logs. Use Jenkins credentials
  (secret text, username/password, SSH keys) injected at runtime with
  \texttt{withCredentials}/credentials bindings, and restrict who can read them.
\end{enumerate}

\section{Conclusion}

The pipeline successfully automated checkout, image build, pytest execution,
deployment on port 8081 and HTTP verification. The failure experiment confirmed that a
failing test blocks deployment and preserves the last working version. The conceptual
Blue--Green discussion shows how the same Jenkins mechanisms can be extended with a
reverse proxy, validation, approval and instant rollback for production releases. As an
optional extension, a GitHub webhook can trigger builds automatically on every push
when the Jenkins server is reachable from GitHub; in this lab, manual builds were used.

\end{document}
